\documentclass[11pt,a4paper]{article}

\usepackage[utf8]{inputenc}
\usepackage[T1]{fontenc}
\usepackage{amsmath,amssymb}
\usepackage{booktabs,tabularx}
\usepackage{enumitem}
\usepackage{geometry}
\usepackage{hyperref}
\usepackage{xcolor}
\usepackage{listings}
\usepackage{fancyhdr}
\usepackage{titlesec}

\geometry{margin=2.2cm,top=2.8cm,bottom=2.8cm}

\definecolor{codebg}{HTML}{F5F5F5}

\hypersetup{
  colorlinks=true,
  linkcolor=blue!60!black,
  urlcolor=blue!60!black,
  pdftitle={ISLS D2 --- Multi-Domain Proof},
  pdfauthor={Sebastian Klemm},
}

\lstset{
  basicstyle=\ttfamily\scriptsize,
  breaklines=true,
  frame=single,
  backgroundcolor=\color{codebg},
  keywordstyle=\color{blue!70!black}\bfseries,
  commentstyle=\color{green!50!black}\itshape,
  numbers=left,
  numberstyle=\tiny\color{gray},
  numbersep=5pt,
  tabsize=4,
  showstringspaces=false,
  xleftmargin=1.5em,
  framexleftmargin=1em,
}

\titleformat{\section}{\Large\bfseries}{\thesection}{1em}{}
\titleformat{\subsection}{\large\bfseries}{\thesubsection}{1em}{}
\titleformat{\subsubsection}{\normalsize\bfseries}{\thesubsubsection}{1em}{}

\pagestyle{fancy}
\fancyhf{}
\fancyhead[L]{\small ISLS D2 --- Multi-Domain Proof}
\fancyhead[R]{\small Klemm, 2026}
\fancyfoot[C]{\thepage}

\begin{document}

% ============================================================
\begin{titlepage}
\begin{center}
\vspace*{1.5cm}
{\Huge\bfseries ISLS D2}\\[0.4cm]
{\LARGE\bfseries Multi-Domain Proof}\\[0.3cm]
{\LARGE\bfseries Resonanz-Breite}\\[1.2cm]
{\Large Drei Dom\"anen. Ein Generator. Drei funktionierende Apps.}\\[1.5cm]
{\large Sebastian Klemm}\\[0.3cm]
{\small 29. M\"arz 2026}\\[2cm]

\begin{tabular}{ll}
\textbf{Repository:} & \texttt{graphity/isls} \\
\textbf{Baut auf:} & D1 (Pipeline $\to$ Compile $\to$ Docker $\to$ Browser) \\
\textbf{Crate:} & \texttt{isls-hypercube} (neue Domains) \\
\textbf{Ziel:} & Beweis: der Generator ist dom\"anenunabh\"angig \\
\end{tabular}

\vfill
{\small
D1 hat bewiesen dass die Pipeline funktioniert: TOML $\to$ HDAG $\to$
kompilierbarer Code $\to$ Docker $\to$ Login $\to$ CRUD.\\[0.2cm]
D2 beweist dass die Pipeline \emph{universell} funktioniert: drei
v\"ollig verschiedene Dom\"anen durchlaufen dieselbe Pipeline und
erzeugen drei funktionierende Anwendungen.\\[0.2cm]
Wenn D2 steht, ist das System kein Warehouse-Generator.\\
Es ist ein App-Generator.
}
\end{center}
\end{titlepage}

\tableofcontents
\newpage

% ============================================================
\section{Ziel}
% ============================================================

Drei TOML-Dateien. Drei v\"ollig verschiedene Dom\"anen. Derselbe Generator.
Alle drei kompilieren, starten in Docker, Login funktioniert, CRUD funktioniert.

\begin{center}
\begin{tabular}{l l l l}
\toprule
\textbf{Dom\"ane} & \textbf{TOML} & \textbf{Entities} & \textbf{Besonderheit} \\
\midrule
Warehouse & \texttt{warehouse.toml} & 7 & D1 --- existiert bereits \\
E-Commerce & \texttt{ecommerce.toml} & 9 & Warenkorb, Bestellstatus, Reviews \\
Projektmanagement & \texttt{project.toml} & 8 & Sprints, Tasks, Kanban-Status \\
\bottomrule
\end{tabular}
\end{center}

% ============================================================
\section{Neue TOML: E-Commerce}
% ============================================================

Datei: \texttt{examples/ecommerce.toml}

\begin{lstlisting}
[app]
name = "ecommerce-platform"
description = "Online store with product catalog, shopping cart, and orders"

[backend]
language = "rust"
framework = "actix-web"
database = "postgresql"
auth_method = "jwt"

[frontend]
type = "spa"
framework = "vanilla"
styling = "minimal"

[deployment]
containerized = true
compose = true

[[entities]]
name = "Category"
fields = [
    { name = "name", type = "String" },
    { name = "slug", type = "String", unique = true },
    { name = "description", type = "String", nullable = true },
    { name = "is_active", type = "bool", default = "true" },
]

[[entities]]
name = "Product"
fields = [
    { name = "name", type = "String" },
    { name = "slug", type = "String", unique = true },
    { name = "description", type = "String", nullable = true },
    { name = "price_cents", type = "i64" },
    { name = "compare_at_price_cents", type = "i64", nullable = true },
    { name = "sku", type = "String", unique = true },
    { name = "inventory_count", type = "i32", default = "0" },
    { name = "is_published", type = "bool", default = "false" },
    { name = "image_url", type = "String", nullable = true },
    { name = "weight_grams", type = "i32", nullable = true },
]
foreign_keys = [
    { target = "Category", nullable = true },
]

[[entities]]
name = "Customer"
fields = [
    { name = "email", type = "String", unique = true },
    { name = "password_hash", type = "String" },
    { name = "first_name", type = "String" },
    { name = "last_name", type = "String" },
    { name = "phone", type = "String", nullable = true },
    { name = "is_active", type = "bool", default = "true" },
]

[[entities]]
name = "Address"
fields = [
    { name = "label", type = "String" },
    { name = "street", type = "String" },
    { name = "city", type = "String" },
    { name = "state", type = "String", nullable = true },
    { name = "postal_code", type = "String" },
    { name = "country", type = "String" },
    { name = "is_default", type = "bool", default = "false" },
]
foreign_keys = [
    { target = "Customer" },
]

[[entities]]
name = "Cart"
fields = [
    { name = "session_id", type = "String" },
    { name = "status", type = "String", default = "active" },
]
foreign_keys = [
    { target = "Customer", nullable = true },
]

[[entities]]
name = "CartItem"
fields = [
    { name = "quantity", type = "i32" },
    { name = "unit_price_cents", type = "i64" },
]
foreign_keys = [
    { target = "Cart" },
    { target = "Product" },
]

[[entities]]
name = "Order"
fields = [
    { name = "order_number", type = "String", unique = true },
    { name = "status", type = "String", default = "pending" },
    { name = "subtotal_cents", type = "i64" },
    { name = "tax_cents", type = "i64", default = "0" },
    { name = "shipping_cents", type = "i64", default = "0" },
    { name = "total_cents", type = "i64" },
    { name = "payment_status", type = "String", default = "unpaid" },
    { name = "notes", type = "String", nullable = true },
]
foreign_keys = [
    { target = "Customer" },
]

[[entities]]
name = "OrderLine"
fields = [
    { name = "quantity", type = "i32" },
    { name = "unit_price_cents", type = "i64" },
]
foreign_keys = [
    { target = "Order" },
    { target = "Product" },
]

[[entities]]
name = "Review"
fields = [
    { name = "rating", type = "i32" },
    { name = "title", type = "String", nullable = true },
    { name = "body", type = "String", nullable = true },
    { name = "is_approved", type = "bool", default = "false" },
]
foreign_keys = [
    { target = "Product" },
    { target = "Customer" },
]

[[entities]]
name = "User"
fields = [
    { name = "email", type = "String", unique = true },
    { name = "password_hash", type = "String" },
    { name = "role", type = "String", default = "user" },
    { name = "is_active", type = "bool", default = "true" },
]
\end{lstlisting}

% ============================================================
\section{Neue TOML: Projektmanagement}
% ============================================================

Datei: \texttt{examples/project.toml}

\begin{lstlisting}
[app]
name = "project-tracker"
description = "Project management with tasks, sprints, and team collaboration"

[backend]
language = "rust"
framework = "actix-web"
database = "postgresql"
auth_method = "jwt"

[frontend]
type = "spa"
framework = "vanilla"
styling = "minimal"

[deployment]
containerized = true
compose = true

[[entities]]
name = "Project"
fields = [
    { name = "name", type = "String" },
    { name = "description", type = "String", nullable = true },
    { name = "status", type = "String", default = "planning" },
    { name = "start_date", type = "String", nullable = true },
    { name = "target_end_date", type = "String", nullable = true },
]

[[entities]]
name = "Sprint"
fields = [
    { name = "name", type = "String" },
    { name = "goal", type = "String", nullable = true },
    { name = "start_date", type = "String" },
    { name = "end_date", type = "String" },
    { name = "status", type = "String", default = "planning" },
]
foreign_keys = [
    { target = "Project" },
]

[[entities]]
name = "Task"
fields = [
    { name = "title", type = "String" },
    { name = "description", type = "String", nullable = true },
    { name = "status", type = "String", default = "todo" },
    { name = "priority", type = "String", default = "medium" },
    { name = "estimate_hours", type = "f64", nullable = true },
    { name = "actual_hours", type = "f64", nullable = true },
    { name = "due_date", type = "String", nullable = true },
    { name = "position", type = "i32", default = "0" },
]
foreign_keys = [
    { target = "Project" },
    { target = "Sprint", nullable = true },
    { target = "User", nullable = true },
]

[[entities]]
name = "Comment"
fields = [
    { name = "body", type = "String" },
]
foreign_keys = [
    { target = "Task" },
    { target = "User" },
]

[[entities]]
name = "Label"
fields = [
    { name = "name", type = "String" },
    { name = "color", type = "String", default = "#3B82F6" },
]
foreign_keys = [
    { target = "Project" },
]

[[entities]]
name = "TaskLabel"
fields = []
foreign_keys = [
    { target = "Task" },
    { target = "Label" },
]

[[entities]]
name = "TeamMember"
fields = [
    { name = "role", type = "String", default = "member" },
]
foreign_keys = [
    { target = "User" },
    { target = "Project" },
]

[[entities]]
name = "User"
fields = [
    { name = "email", type = "String", unique = true },
    { name = "password_hash", type = "String" },
    { name = "display_name", type = "String" },
    { name = "avatar_url", type = "String", nullable = true },
    { name = "role", type = "String", default = "user" },
    { name = "is_active", type = "bool", default = "true" },
]
\end{lstlisting}

% ============================================================
\section{Was sich \"andern muss}
% ============================================================

\subsection{TOML-Parser in \texttt{isls-hypercube}}

Die TOML-Dateien oben verwenden ein \textbf{generisches Format} ---
Entities werden direkt in der TOML definiert, nicht \"uber hardcodierte
Domain-Templates. Der Parser muss:

\begin{enumerate}[nosep]
\item \texttt{[[entities]]} Array parsen: Name, Fields, Foreign Keys.
\item Jedes Field hat: \texttt{name}, \texttt{type}, optional \texttt{nullable},
      \texttt{unique}, \texttt{default}.
\item Jede Foreign Key hat: \texttt{target} (Entity-Name), optional \texttt{nullable}.
\item Das Ergebnis ist ein \texttt{AppSpec} mit \texttt{Vec<EntityDef>}.
\end{enumerate}

\textbf{KRITISCH:} Wenn der aktuelle Parser Entity-Definitionen aus
hardcodierten Domain-Templates l\"adt (z.B. \texttt{warehouse\_domain()}),
muss er auf das TOML-basierte Format umgestellt werden. Die TOML ist
die Wahrheit. Keine hardcodierten Domains mehr.

Wenn der Parser bereits das generische Format unterst\"utzt, reicht es
die zwei neuen TOML-Dateien anzulegen.

\subsection{Pr\"ufe: \texttt{EntityDef} Struktur}

Die \texttt{EntityDef} in \texttt{isls-types} muss unterst\"utzen:

\begin{lstlisting}[language=C]
pub struct EntityDef {
    pub name: String,              // "Product"
    pub fields: Vec<FieldDef>,
    pub foreign_keys: Vec<ForeignKeyDef>,
}

pub struct FieldDef {
    pub name: String,              // "price_cents"
    pub field_type: String,        // "i64"
    pub nullable: bool,            // false
    pub unique: bool,              // false
    pub default: Option<String>,   // Some("0")
}

pub struct ForeignKeyDef {
    pub target: String,            // "Category"
    pub nullable: bool,            // true
}
\end{lstlisting}

Wenn \texttt{FieldDef} noch kein \texttt{unique} oder \texttt{default}
hat, hinzuf\"ugen. Wenn \texttt{ForeignKeyDef} nicht existiert, erstellen.

\subsection{Strukturelle Generatoren: FK-Felder}

Die strukturellen Generatoren m\"ussen Foreign Keys korrekt verarbeiten:

\begin{itemize}[nosep]
\item \textbf{Model:} F\"ur jede FK ein Feld \texttt{\{target\_snake\}\_id: i64}
      (oder \texttt{Option<i64>} wenn nullable).
\item \textbf{Migration:} F\"ur jede FK ein
      \texttt{\{target\_snake\}\_id BIGINT [NOT NULL] REFERENCES \{target\_table\}(id)}.
      Tabellen m\"ussen in topologischer Reihenfolge erzeugt werden
      (referenzierte Tabelle vor referenzierender).
\item \textbf{CreatePayload:} FK-Felder als \texttt{i64} (required) oder
      \texttt{Option<i64>} (nullable FK).
\item \textbf{UpdatePayload:} FK-Felder immer als \texttt{Option<i64>}.
\item \textbf{Query INSERT:} FK-Felder im INSERT-Statement.
\end{itemize}

\subsection{Migration: Topologische Tabellenreihenfolge}

Die Migration muss Tabellen in der richtigen Reihenfolge erzeugen.
\texttt{products} referenziert \texttt{categories} --- also muss
\texttt{categories} zuerst kommen. \texttt{users} hat keine FKs ---
kommt ganz am Anfang.

\begin{lstlisting}[language=C]
fn topological_sort_entities(entities: &[EntityDef]) -> Vec<&EntityDef> {
    // Entities ohne FKs zuerst (users, categories, projects, labels)
    // Dann Entities deren FKs bereits erzeugt wurden
    // Zyklus-Erkennung: wenn A -> B -> A, Fehler
}
\end{lstlisting}

\subsection{Plural-Namen}

Der Generator braucht korrekte Pluralformen:

\begin{center}
\small
\begin{tabular}{l l}
\toprule
\textbf{Entity} & \textbf{Plural} \\
\midrule
Product & products \\
Category & categories \\
Address & addresses \\
OrderLine & order\_lines \\
CartItem & cart\_items \\
StockMovement & stock\_movements \\
TeamMember & team\_members \\
TaskLabel & task\_labels \\
\bottomrule
\end{tabular}
\end{center}

Einfache Regel: Wenn der snake\_case-Name auf \texttt{s} endet,
kein zus\"atzliches \texttt{s}. Wenn auf \texttt{y} endet (nach
Konsonant), ersetze mit \texttt{ies}. Wenn auf \texttt{ss}/\texttt{sh}/\texttt{ch}/\texttt{x}
endet, h\"ange \texttt{es} an. Sonst h\"ange \texttt{s} an.

Oder: ein \texttt{plural\_name} Feld in EntityDef das optional
\"uberschrieben werden kann.

\subsection{Entity mit 0 eigenen Feldern}

\texttt{TaskLabel} hat keine eigenen Felder --- nur zwei FKs
(\texttt{task\_id}, \texttt{label\_id}). Das ist eine reine
Join-Tabelle. Die Generatoren m\"ussen damit umgehen:

\begin{itemize}[nosep]
\item Model: Struct hat nur \texttt{id}, \texttt{task\_id},
      \texttt{label\_id}, \texttt{created\_at}, \texttt{updated\_at}.
\item CreatePayload: nur \texttt{task\_id: i64}, \texttt{label\_id: i64}.
\item UpdatePayload: nur \texttt{task\_id: Option<i64>},
      \texttt{label\_id: Option<i64>}.
\end{itemize}

% ============================================================
\section{Verifikation}
% ============================================================

\begin{lstlisting}[language=bash]
# 1. Workspace
cargo build --workspace --release
cargo test --workspace

# 2. Warehouse (existiert, Regression-Test)
rm -rf output/warehouse
cargo run --release -p isls-cli -- forge-v2 \
    --requirements examples/warehouse.toml \
    --api-key %OPENAI_API_KEY% --output ./output/warehouse
cd output/warehouse/backend && cargo build && cd ../../..

# 3. E-Commerce (NEU)
rm -rf output/ecommerce
cargo run --release -p isls-cli -- forge-v2 \
    --requirements examples/ecommerce.toml \
    --api-key %OPENAI_API_KEY% --output ./output/ecommerce
cd output/ecommerce/backend && cargo build && cd ../../..

# 4. Projektmanagement (NEU)
rm -rf output/project
cargo run --release -p isls-cli -- forge-v2 \
    --requirements examples/project.toml \
    --api-key %OPENAI_API_KEY% --output ./output/project
cd output/project/backend && cargo build && cd ../../..

# 5. Docker-Test fuer alle drei
for dir in warehouse ecommerce project; do
    cd output/$dir
    docker-compose up -d --build
    timeout 30
    curl -s http://localhost:8080/api/health
    curl -s -X POST http://localhost:8080/api/auth/login \
        -H "Content-Type: application/json" \
        -d "{\"email\":\"admin@example.com\",\"password\":\"admin123\"}"
    docker-compose down -v
    cd ../..
done

# Alle drei muessen:
# - cargo build: 0 errors
# - health: 200 OK
# - login: 200 OK mit JWT
\end{lstlisting}

% ============================================================
\section{Constraints f\"ur Claude Code}
% ============================================================

\begin{enumerate}
\item \textbf{Keine hardcodierten Domain-Templates.} Entities kommen
      aus der TOML, nicht aus \texttt{warehouse\_domain()} oder
      \texttt{ecommerce\_domain()}. Die TOML ist die einzige Wahrheit.
\item \textbf{Alle drei TOMLs m\"ussen den identischen Pipeline-Pfad
      durchlaufen.} Kein Spezialcode f\"ur einzelne Domains.
\item \textbf{Foreign Keys m\"ussen funktionieren.} Model, Migration,
      CreatePayload, Queries --- alles muss FK-Felder korrekt
      verarbeiten.
\item \textbf{Migration: topologische Reihenfolge.} Referenzierte
      Tabellen zuerst.
\item \textbf{Join-Tabellen (0 eigene Felder) m\"ussen funktionieren.}
\item \textbf{D1 darf nicht brechen.} \texttt{warehouse.toml} muss
      weiterhin kompilierbaren Output erzeugen.
\item \textbf{Die TOML-Dateien in dieser Spec sind exakt.} Nicht
      \"andern. Die Entity-Definitionen, Felder, Foreign Keys und
      Defaults sind absichtlich so gew\"ahlt.
\item \textbf{Alle bestehenden Tests m\"ussen bestehen.}
\item \textbf{Ein Commit nach Abschluss.}
\item \textbf{Mock-Modus und LLM-Modus m\"ussen f\"ur alle drei
      Domains kompilierbaren Output erzeugen.}
\end{enumerate}

% ============================================================
\section{Was D2 beweist}
% ============================================================

Wenn alle drei Domains kompilieren, starten und CRUD funktioniert:

\begin{itemize}[nosep]
\item Der Generator ist dom\"anenunabh\"angig.
\item Die strukturellen Generatoren arbeiten korrekt mit beliebigen
      EntityDefs, nicht nur mit Warehouse-Entities.
\item Foreign Keys, Join-Tabellen und verschiedene Feldtypen
      werden korrekt verarbeitet.
\item Die LLM-generierten SQL-Queries funktionieren f\"ur
      verschiedene Schemata.
\item Der HDAG skaliert: 7 Entities (Warehouse), 9 Entities
      (E-Commerce), 8 Entities (Projekt).
\end{itemize}

D2 ist der Beweis dass das System kein Warehouse-Drucker ist.
Es ist ein App-Drucker.

\vfill
\noindent\rule{\textwidth}{0.4pt}
\begin{center}
\small
ISLS D2 --- Multi-Domain Proof --- Sebastian Klemm --- 29. M\"arz 2026\\
Drei Dom\"anen. Ein Generator. Drei funktionierende Apps.\\
Warehouse. E-Commerce. Projektmanagement.\\
Derselbe HDAG. Derselbe Pipeline. Drei verschiedene Welten.\\
\url{https://github.com/LashSesh/graphity}
\end{center}

\end{document}
